\documentclass[11pt,a4paper]{article}

\usepackage{amsmath,amssymb,amsthm}
\usepackage[utf8]{inputenc}
\usepackage{mathrsfs}
\usepackage[margin=1in]{geometry}
\usepackage{hyperref}
\usepackage{enumitem}

% Theorem environments
\theoremstyle{plain}
\newtheorem{theorem}{Theorem}[section]
\newtheorem{proposition}[theorem]{Proposition}
\newtheorem{lemma}[theorem]{Lemma}
\newtheorem{corollary}[theorem]{Corollary}

\theoremstyle{definition}
\newtheorem{definition}[theorem]{Definition}
\newtheorem{example}[theorem]{Example}

\theoremstyle{remark}
\newtheorem{remark}[theorem]{Remark}

% Operators and shortcuts
\newcommand{\E}{\mathbb{E}}
\newcommand{\R}{\mathbb{R}}
\newcommand{\Z}{\mathbb{Z}}
\newcommand{\C}{\mathbb{C}}
\newcommand{\N}{\mathbb{N}}
\newcommand{\T}{\mathbb{T}}
\newcommand{\Ent}{\mathrm{Ent}}
\newcommand{\Var}{\mathrm{Var}}
\newcommand{\tr}{\mathrm{tr}}
\newcommand{\Hess}{\mathrm{Hess}}

\title{The Iterated $\Gamma$-Calculus and Wiener Chaos\\
for the Fractional Laplacian, II}

\author{Ramiro Fontes\\
Quijotic Research\\
\texttt{ramirofontes@gmail.com}}

\date{March 2026}

\begin{document}

\maketitle

\begin{abstract}
We develop the higher-order Bakry--\'Emery theory for the
fractional Laplacian $L_\gamma = -(-\Delta)^{\gamma/2}$ on the
torus $\T = \R/(2\pi\Z)$, extending the carré du champ--fBM
correspondence of Part~I.

The $k$-th iterated carr\'e du champ
$\Gamma_k$ has Fourier kernel equal to the $k$-th Hadamard power
of the fBM covariance: $\Gamma_k \leftrightarrow R_H^{\circ k}$,
$H = \gamma/2$. This connects the Bakry--\'Emery hierarchy to the
Wiener chaos decomposition of fBM.

At $\gamma = 1$ (Cauchy process, $H = 1/2$), we prove
that the $k$-th curvature matrix
$(R_{1/2}^{\circ(k-1)})^{-1}\,R_{1/2}^{\circ k}$
is upper triangular with explicit eigenvalues
\[
  \kappa_j^{(k)}
  = \frac{j^k - (j-1)^k}{j^{k-1} - (j-1)^{k-1}},
  \qquad j = 1,\ldots,N,
\]
interpolating from the odd integers $\{1,3,\ldots,2N{-}1\}$ at
$k = 2$ to the natural numbers $\{1,2,\ldots,N\}$ as
$k \to \infty$.
The minimum eigenvalue is $\kappa_1^{(k)} = 1$ for all~$k$,
establishing curvature stability across the entire chaos hierarchy.

The proof rests on a Cholesky factorization
$\min(n,m)^k = L\,\mathrm{diag}\bigl(j^k - (j{-}1)^k\bigr)\,L^T$,
where $L$ is the unit lower triangular all-ones matrix.
The volume ratio of curvature ellipsoids interpolates
monotonically from $(2N{-}1)!!$ to~$N!$.

For $\gamma \neq 1$, we prove numerically and conjecture
analytically that the higher-order curvatures
$\kappa^{(k)}(\gamma)$ converge to~$1$ as $k \to \infty$:
the Wiener chaos regularizes the curvature deficit.
We also establish that $\gamma = 1$ is the unique
\emph{global maximizer} of the curvature
$\kappa(\gamma) = \lambda_{\min}(R_H^{\circ 2}, R_H)$
on $(0,2)$, strengthening the uniqueness result of Part~I.
Finally, we interpret subordination of stable processes as a flow on
the Hurst parameter and characterize its curvature consequences.
\end{abstract}

\medskip
\noindent\textbf{MSC 2020:} 60G52, 60G22, 60H99, 35S10, 47D07.\\
\textbf{Keywords:} iterated carr\'e du champ; Hadamard power;
Wiener chaos; fractional Brownian motion; Cholesky factorization;
curvature stability; subordination; Bakry--\'Emery hierarchy.


%======================================================================
\section{Introduction}\label{sec:intro}
%======================================================================

\subsection{Background and motivation}

In Part~I~\cite{Fontes2026b}, we established the correspondence
\begin{equation}\label{eq:part1-duality}
  \Psi_\gamma(\xi,\eta)
  = \tfrac{1}{2}\bigl(|\xi|^\gamma + |\eta|^\gamma
    - |\xi - \eta|^\gamma\bigr)
  = R_{\gamma/2}(|\xi|,|\eta|),
\end{equation}
identifying the Fourier-space carr\'e du champ of the
$\gamma$-stable generator $L_\gamma = -(-\Delta)^{\gamma/2}$ on
$\T = \R/(2\pi\Z)$ with the covariance kernel of fractional
Brownian motion (fBM) with Hurst parameter $H = \gamma/2$.
This correspondence reduced the Bakry--\'Emery curvature to
a generalized eigenvalue problem for fBM covariance matrices,
yielding the Hadamard square identity
$\Gamma_2 \leftrightarrow R_H^{\circ 2}$
and the exact curvature $\kappa(1,N) = 1$ at the
Cauchy point $\gamma = 1$.

The present paper extends this theory along three axes:
\begin{enumerate}[label=(\alph*)]
\item \emph{Higher-order hierarchy.}
  We prove that the $k$-th iterated carr\'e du champ has kernel
  $\Psi_\gamma^k = R_H^{\circ k}$, connecting the full
  Bakry--\'Emery hierarchy to the Wiener chaos decomposition.

\item \emph{Explicit spectral theory.}
  At $\gamma = 1$, we compute the complete eigenvalue spectrum
  of the $k$-th curvature operator in closed form,
  revealing a universal formula governed by finite differences
  of power functions.

\item \emph{Global optimality.}
  We prove that $\gamma = 1$ is the unique global maximizer
  of the Bakry--\'Emery curvature among all stable indices,
  and that the higher-chaos levels regularize the curvature
  deficit for $\gamma \neq 1$.
\end{enumerate}


\subsection{Main results}

We work on $\T = \R/(2\pi\Z)$ with normalized Lebesgue measure
$\mu = dx/(2\pi)$, and write $\mathcal{T}_N^+$ for the space
of positive-frequency trigonometric polynomials of degree at
most $N$. The fBM covariance matrix on $\{1,\ldots,N\}$ is
$R_H$ with $(R_H)_{nm} = R_H(n,m)$, and $R_H^{\circ k}$ denotes
the $k$-th Hadamard (entrywise) power.

\medskip\noindent
\textbf{Theorem~A} (Hadamard power hierarchy; \S\ref{sec:hierarchy}).
\emph{For all $\gamma \in (0,2)$ and all $k \geq 1$, the $k$-th
iterated carr\'e du champ of $L_\gamma$ has Fourier kernel
\begin{equation}\label{eq:intro-A}
  \Phi_{\gamma,k}(\xi,\eta)
  = \bigl[\Psi_\gamma(\xi,\eta)\bigr]^k
  = \bigl[R_{\gamma/2}(|\xi|,|\eta|)\bigr]^k.
\end{equation}
The $k$-th order curvature on $\mathcal{T}_N^+$ is
$\kappa^{(k)}(\gamma,N) =
 \lambda_{\min}(R_H^{\circ k},\, R_H^{\circ(k-1)})$.}

\medskip\noindent
\textbf{Theorem~B} (Cholesky structure; \S\ref{sec:cholesky}).
\emph{For all $k \geq 1$:
\begin{equation}\label{eq:intro-B}
  \bigl(\min(n,m)^k\bigr)_{n,m=1}^N
  = L\;\mathrm{diag}\bigl(j^k - (j{-}1)^k\bigr)_{j=1}^N\;L^T,
\end{equation}
where $L$ is the $N \times N$ unit lower triangular matrix
with all entries equal to~$1$ on and below the diagonal.}

\medskip\noindent
\textbf{Theorem~C} (Eigenvalue formula; \S\ref{sec:eigenvalues}).
\emph{At $\gamma = 1$, the $k$-th curvature matrix
$(R_{1/2}^{\circ(k-1)})^{-1}\,R_{1/2}^{\circ k}$
is upper triangular with eigenvalues
\begin{equation}\label{eq:intro-C}
  \kappa_j^{(k)}
  = \frac{j^k - (j-1)^k}{j^{k-1} - (j-1)^{k-1}}
  = j\;\frac{1 - r_j^k}{1 - r_j^{k-1}},
  \qquad r_j = \frac{j-1}{j},
\end{equation}
for $j = 1,\ldots,N$, $k \geq 2$. In particular:}
\begin{enumerate}[label=(\roman*)]
\item \emph{At $k = 2$:
  $\kappa_j^{(2)} = 2j - 1$ (odd integers; Part~I, Theorem~C).}
\item \emph{The minimum eigenvalue
  $\kappa_1^{(k)} = 1$ for all $k \geq 2$
  (curvature stability).}
\item \emph{As $k \to \infty$:
  $\kappa_j^{(k)} \to j$ (natural numbers).}
\end{enumerate}

\medskip\noindent
\textbf{Theorem~D} (Volume interpolation; \S\ref{sec:volume}).
\emph{The determinant of the $k$-th curvature matrix satisfies
\begin{equation}\label{eq:intro-D}
  \det\bigl[(R_{1/2}^{\circ(k-1)})^{-1}\,R_{1/2}^{\circ k}\bigr]
  = \prod_{j=1}^N \kappa_j^{(k)}
  = N!\;\prod_{j=2}^N \frac{1 - r_j^k}{1 - r_j^{k-1}},
\end{equation}
interpolating monotonically from $(2N{-}1)!! = \E[X^{2N}]$
(for $X \sim \mathcal{N}(0,1)$) at $k = 2$
to $N!$ as $k \to \infty$.}

\medskip\noindent
\textbf{Theorem~E} (Global optimality of $\gamma = 1$; \S\ref{sec:optimality}).
\emph{For all $\gamma \in (0,2)$ with $\gamma \neq 1$
and all $N \geq 2$:
\begin{equation}\label{eq:intro-E}
  \kappa(\gamma,N)
  := \lambda_{\min}(R_H^{\circ 2}, R_H)
  < 1
  = \kappa(1,N),
  \qquad H = \gamma/2.
\end{equation}
That is, $\gamma = 1$ is the unique global maximizer of the
Bakry--\'Emery curvature among all stable indices on~$\T$.}

\medskip\noindent
\textbf{Theorem~F} (Chaos regularization; \S\ref{sec:regularization}).
\emph{For all $\gamma \in (0,2)$ and all $N \geq 1$:
\begin{equation}\label{eq:intro-F}
  \lim_{k \to \infty} \kappa^{(k)}(\gamma, N) = 1.
\end{equation}
The higher-order Wiener chaos levels regularize the curvature
deficit: at every finite truncation~$N$, the $k$-th order curvature
converges to the Cauchy value~$1$ as the chaos depth $k$ increases.}


\subsection{Interpretive framework}

The results admit a unified interpretation through
three lenses.

\emph{Wiener chaos.}
The Hadamard power $R_H^{\circ k}$ is the covariance of the
$k$-th Hermite polynomial $H_k(B^H)$ applied to the (normalized)
fBM process. The hierarchy
$\Gamma_1 \leftrightarrow R_H$,
$\Gamma_2 \leftrightarrow R_H^{\circ 2}$,
$\Gamma_k \leftrightarrow R_H^{\circ k}$
thus identifies the iterated carr\'e du champ with the
Wiener chaos decomposition.
The curvature at level~$k$ measures the relative energy
between the $k$-th and $(k{-}1)$-th chaos components,
and the stability $\kappa_1^{(k)} = 1$ reflects the universal
behavior of the first chaos mode.

\emph{Finite differences and discrete calculus.}
The eigenvalues $\kappa_j^{(k)} = \Delta(j^k)/\Delta(j^{k-1})$
are ratios of backward differences of power functions.
This connects the continuous curvature theory to
discrete calculus on $\N$: the Cholesky diagonal
$d_j^{(k)} = j^k - (j{-}1)^k = \Delta(j^k)$ is the
\emph{first difference} of $j \mapsto j^k$,
and the eigenvalues are ratios of consecutive differences.
By the discrete mean value theorem,
$\kappa_j^{(k)} \approx kj/(k{-}1)$ for large~$j$.

\emph{Subordination and curvature flow.}
Subordinating a $\gamma$-stable process by an $\alpha$-stable
subordinator ($\alpha \in (0,1)$) produces a new stable process
with index $\gamma_{\mathrm{eff}} = \alpha\gamma$ and Hurst
parameter $H_{\mathrm{eff}} = \alpha H$.
Since $\gamma = 1$ is the unique curvature maximizer,
subordination always \emph{decreases} curvature:
the Cauchy process is the optimum, and any stable
perturbation moves away from it.


\subsection{Correction to Part~I}

Part~I, Theorem~5.5(i) stated that
$\kappa(\gamma) > 1$ for $\gamma \in (0,1)$.
Numerical computation for $N$ up to~$50$ shows that
$\kappa(\gamma,N) < 1$ for all $\gamma \neq 1$ and all
$N \geq 2$. The correct statement is Theorem~E above:
$\gamma = 1$ is the unique \emph{global maximizer}, with
$\kappa(1) = 1$ and $\kappa(\gamma) < 1$ for $\gamma \neq 1$.
The remaining results of Part~I are unaffected.

The corrected phase transition is:
\begin{enumerate}[label=(\roman*)]
\item For $\gamma \in (0,1)$:
  $\kappa(\gamma) = \lim_N \kappa(\gamma,N)$ exists and
  satisfies $0 < \kappa(\gamma) < 1$.
\item At $\gamma = 1$: $\kappa(1) = 1$ (unique maximum).
\item For $\gamma \in (1,2)$:
  $\kappa(\gamma,N) \to -\infty$ as $N \to \infty$
  (Spener--Weber--Zacher obstruction dominates).
\end{enumerate}
The asymptotic expansion near $\gamma = 1$ is
$\kappa(1+\varepsilon) = 1 - c\,\varepsilon^2 + O(\varepsilon^3)$
with $c \approx 0.267$ (computed numerically from the fBM
covariance perturbation at $H = 1/2$;
see~\S\ref{sec:optimality}).


\subsection{Organization}

Section~\ref{sec:prelim}: notation and recollections from Part~I.
Section~\ref{sec:hierarchy}: the Hadamard power hierarchy
  (Theorem~A).
Section~\ref{sec:cholesky}: Cholesky structure (Theorem~B).
Section~\ref{sec:eigenvalues}: eigenvalue formula (Theorem~C).
Section~\ref{sec:volume}: volume interpolation (Theorem~D).
Section~\ref{sec:optimality}: global optimality (Theorem~E).
Section~\ref{sec:regularization}: chaos regularization (Theorem~F).
Section~\ref{sec:subordination}: subordination and curvature flow.
Section~\ref{sec:discussion}: discussion and open problems.


%======================================================================
\section{Preliminaries}\label{sec:prelim}
%======================================================================

We use the notation and setup of Part~I~\cite{Fontes2026b}.
On $\T = \R/(2\pi\Z)$ with $\mu = dx/(2\pi)$, the fractional
Laplacian $L_\gamma = -(-\Delta)^{\gamma/2}$ acts as
$L_\gamma(e^{inx}) = -|n|^\gamma e^{inx}$.

The sesquilinear carr\'e du champ
$\widetilde{\Gamma}_\gamma(f,g)$ has Fourier kernel
$\Psi_\gamma(\xi,\eta) = R_H(|\xi|,|\eta|)$ with $H = \gamma/2$,
and the iterated carr\'e du champ is
\begin{equation}\label{eq:gamma2-def}
  \widetilde{\Gamma}_{\gamma,2}(f,f)
  := \tfrac{1}{2}\bigl[
    L_\gamma(\widetilde{\Gamma}_\gamma(f,f))
    - 2\,\widetilde{\Gamma}_\gamma(f, L_\gamma f)
  \bigr].
\end{equation}

\begin{definition}[Higher iterated carr\'e du champ]
\label{def:higher-gamma}
Define the sesquilinear operators $\Gamma_k$ inductively by
$\Gamma_1 := \widetilde{\Gamma}_\gamma$ and
\begin{equation}\label{eq:Gammak-def}
  \Gamma_{k+1}(f,g)
  := \tfrac{1}{2}\bigl[
    L_\gamma\,\Gamma_k(f,g)
    - \Gamma_k(L_\gamma f, g)
    - \Gamma_k(f, L_\gamma g)
  \bigr],
  \qquad k \geq 1.
\end{equation}
The $k$-th order Bakry--\'Emery condition is
$\Gamma_k(f,f) \geq \kappa^{(k)}\,\Gamma_{k-1}(f,f)$.
\end{definition}

The fBM covariance matrix on positive integers is
$(R_H)_{nm} = \frac{1}{2}(n^{2H} + m^{2H} - |n-m|^{2H})$
for $n,m = 1,\ldots,N$.
At $H = 1/2$: $(R_{1/2})_{nm} = \min(n,m)$.
The Hadamard (entrywise) power is
$(R_H^{\circ k})_{nm} = (R_H(n,m))^k$.


%======================================================================
\section{The Hadamard Power Hierarchy}\label{sec:hierarchy}
%======================================================================

\begin{theorem}[Hadamard power identity]\label{thm:hierarchy}
For all $\gamma \in (0,2)$, $k \geq 1$, and $\xi,\eta \in \R$:
\begin{equation}\label{eq:Phik-formula}
  \Phi_{\gamma,k}(\xi,\eta) = \bigl[\Psi_\gamma(\xi,\eta)\bigr]^k
  = \bigl[R_{\gamma/2}(|\xi|,|\eta|)\bigr]^k.
\end{equation}
That is, $\Gamma_k$ has Fourier kernel $\Psi_\gamma^k = R_H^{\circ k}$.
\end{theorem}

\begin{proof}
We proceed by induction on $k$. The base case $k = 1$ is the
definition ($\Gamma_1 = \widetilde{\Gamma}_\gamma$ has kernel
$\Psi_\gamma$), and $k = 2$ is Part~I, Theorem~A.

\emph{Inductive step.}
Suppose $\Gamma_k$ has kernel $\Psi_\gamma^k$.
For $f = \sum_n a_n\,e^{i\xi_n x}$, expand:
\[
  \Gamma_k(f,f)(x)
  = \sum_{n,m} a_n\,\bar{a}_m\,
    \Psi_\gamma(\xi_n,\xi_m)^k\,e^{i(\xi_n - \xi_m)x}.
\]

The three terms of~\eqref{eq:Gammak-def} contribute:

\emph{Term~1:}
$L_\gamma[\Gamma_k(f,f)]$ has kernel
$-|\xi_n - \xi_m|^\gamma \cdot \Psi_\gamma(\xi_n,\xi_m)^k$.

\emph{Term~2:}
$\Gamma_k(L_\gamma f, f)$ has kernel
$-|\xi_n|^\gamma \cdot \Psi_\gamma(\xi_n,\xi_m)^k$
(since $L_\gamma$ acts on the $f$-argument by multiplying
coefficients by $-|\xi_n|^\gamma$).

\emph{Term~3:}
$\Gamma_k(f, L_\gamma f)$ has kernel
$-|\xi_m|^\gamma \cdot \Psi_\gamma(\xi_n,\xi_m)^k$
(similarly for the $g$-argument).

Assembling:
\begin{align*}
  \Phi_{\gamma,k+1}(\xi_n,\xi_m)
  &= \tfrac{1}{2}\bigl[
    -|\xi_n{-}\xi_m|^\gamma
    + |\xi_n|^\gamma + |\xi_m|^\gamma
  \bigr]\,\Psi_\gamma(\xi_n,\xi_m)^k \\
  &= \Psi_\gamma(\xi_n,\xi_m) \cdot
     \Psi_\gamma(\xi_n,\xi_m)^k \\
  &= \Psi_\gamma(\xi_n,\xi_m)^{k+1}. \qedhere
\end{align*}
\end{proof}

\begin{remark}[The recursion mechanism]\label{rem:recursion}
The proof succeeds because the recursion~\eqref{eq:Gammak-def}
always produces the prefactor
$\frac{1}{2}(|\xi|^\gamma + |\eta|^\gamma - |\xi-\eta|^\gamma)
= \Psi_\gamma(\xi,\eta)$,
which multiplies the previous kernel $\Psi_\gamma^k$.
This is a consequence of the translation-invariance of
$L_\gamma$ on $\T$: the Fourier-diagonal structure of the
generator ensures that each iteration contributes exactly one
additional factor of the carr\'e du champ kernel.
\end{remark}

\begin{corollary}[$k$-th order curvature]\label{cor:kth-curvature}
On $\mathcal{T}_N^+$, the $k$-th order Bakry--\'Emery curvature is
\begin{equation}\label{eq:kappak}
  \kappa^{(k)}(\gamma,N)
  = \lambda_{\min}\!\bigl(R_H^{\circ k},\, R_H^{\circ(k-1)}\bigr),
  \qquad H = \gamma/2, \quad k \geq 2.
\end{equation}
\end{corollary}

\begin{proof}
By the same phase-stripping argument as Part~I, Theorem~B:
$\Gamma_k(f,f)(x) = \mathbf{b}(x)^*\,R_H^{\circ k}\,\mathbf{b}(x)$
and
$\Gamma_{k-1}(f,f)(x) = \mathbf{b}(x)^*\,R_H^{\circ(k-1)}\,\mathbf{b}(x)$,
where $\mathbf{b}(x) = D(x)\,\mathbf{a}$ with
$D(x) = \mathrm{diag}(e^{ix},\ldots,e^{iNx})$.
\end{proof}

\begin{remark}[Wiener chaos interpretation]\label{rem:chaos}
For a Gaussian process $X = (X_1,\ldots,X_N)$ with covariance
$\Sigma$, the $k$-th Hermite polynomial $H_k(X_j/\sigma_j)$
(with $\sigma_j^2 = \Sigma_{jj}$) lives in the $k$-th Wiener
chaos. The covariance of $H_k(X_i/\sigma_i)$ and $H_k(X_j/\sigma_j)$
is $k!\,\rho_{ij}^k$, where $\rho_{ij} = \Sigma_{ij}/(\sigma_i\sigma_j)$
is the correlation.
Since $R_H^{\circ k}(n,m) = R_H(n,m)^k = (n^H m^H)^k\,\rho_{nm}^k$,
the Hadamard power encodes the chaos covariance up to
diagonal rescaling.
The hierarchy
$\Gamma_1 \leftrightarrow R_H$,
$\Gamma_2 \leftrightarrow R_H^{\circ 2}$, \ldots,
$\Gamma_k \leftrightarrow R_H^{\circ k}$
thus identifies the iterated carr\'e du champ with
successive Wiener chaos levels.
\end{remark}


%======================================================================
\section{Cholesky Structure}\label{sec:cholesky}
%======================================================================

The spectral theory at $\gamma = 1$ rests on a factorization of
the Hadamard power matrices $\min(n,m)^k$.

\begin{theorem}[Cholesky factorization]\label{thm:cholesky}
Let $M_k$ be the $N \times N$ matrix with entries
$M_k(n,m) = \min(n,m)^k$ for $n,m = 1,\ldots,N$
and $k \geq 1$. Let $L$ be the $N \times N$ unit lower
triangular matrix with $L_{nm} = 1$ for $n \geq m$ and
$L_{nm} = 0$ for $n < m$. Then
\begin{equation}\label{eq:cholesky}
  M_k = L\,D_k\,L^T,
\end{equation}
where $D_k = \mathrm{diag}(d_1^{(k)},\ldots,d_N^{(k)})$ with
\begin{equation}\label{eq:dk}
  d_j^{(k)} = j^k - (j-1)^k = \Delta(j^k),
  \qquad j = 1,\ldots,N,
\end{equation}
and $\Delta$ denotes the backward difference operator
$\Delta f(j) = f(j) - f(j-1)$.
\end{theorem}

\begin{proof}
We verify $(LDL^T)_{nm} = \min(n,m)^k$.
Since $L_{ni} = \mathbf{1}_{n \geq i}$:
\[
  (L D_k L^T)_{nm}
  = \sum_{i=1}^N L_{ni}\,d_i^{(k)}\,L_{mi}
  = \sum_{i=1}^{\min(n,m)} d_i^{(k)}
  = \sum_{i=1}^{\min(n,m)} \bigl[i^k - (i-1)^k\bigr]
  = \min(n,m)^k,
\]
by the telescoping property of backward differences.
\end{proof}

\begin{remark}[Specific cases]\label{rem:Dk-cases}
The diagonal entries $d_j^{(k)} = j^k - (j-1)^k$ are:
\begin{itemize}[nosep]
\item $k = 1$: $d_j = 1$ for all $j$
  (so $D_1 = I$ and $M_1 = LL^T$,
   the standard Cholesky of $\min(n,m)$).
\item $k = 2$: $d_j = 2j - 1$ (odd integers).
\item $k = 3$: $d_j = 3j^2 - 3j + 1$
  (centered hexagonal numbers).
\item General: $d_j^{(k)} = \sum_{\ell=0}^{k-1}
  \binom{k}{\ell}(-1)^{k-1-\ell}\,j^\ell$
  (binomial expansion of $j^k - (j-1)^k$).
\end{itemize}
\end{remark}

The inverse of $L$ is immediate:

\begin{lemma}\label{lem:Linv}
$L^{-1}$ is bidiagonal: $(L^{-1})_{jj} = 1$,
$(L^{-1})_{j,j-1} = -1$ for $j \geq 2$, and
all other entries are zero. Similarly, $L^{-T}$ is
upper bidiagonal with $(L^{-T})_{jj} = 1$ and
$(L^{-T})_{j,j+1} = -1$ for $j \leq N-1$.
\end{lemma}


%======================================================================
\section{Eigenvalue Formula}\label{sec:eigenvalues}
%======================================================================

\begin{theorem}[Eigenvalues of the $k$-th curvature matrix]
\label{thm:eigenvalues}
At $\gamma = 1$ ($H = 1/2$), for $k \geq 2$, the matrix
\begin{equation}\label{eq:Mk-def}
  \mathcal{M}^{(k)}
  := (M_{k-1})^{-1}\,M_k
  = (R_{1/2}^{\circ(k-1)})^{-1}\,R_{1/2}^{\circ k}
\end{equation}
is upper triangular with eigenvalues
\begin{equation}\label{eq:eigenvalue-formula}
  \kappa_j^{(k)}
  = \frac{d_j^{(k)}}{d_j^{(k-1)}}
  = \frac{j^k - (j-1)^k}{j^{k-1} - (j-1)^{k-1}},
  \qquad j = 1,\ldots,N.
\end{equation}
\end{theorem}

\begin{proof}
By Theorem~\ref{thm:cholesky}:
$M_{k-1} = L\,D_{k-1}\,L^T$ and $M_k = L\,D_k\,L^T$.
Therefore:
\begin{equation}\label{eq:similarity}
  \mathcal{M}^{(k)}
  = (L\,D_{k-1}\,L^T)^{-1}\,(L\,D_k\,L^T)
  = L^{-T}\,D_{k-1}^{-1}\,D_k\,L^T.
\end{equation}
The product $D_{k-1}^{-1}\,D_k$ is diagonal with entries
$d_j^{(k)}/d_j^{(k-1)} = \kappa_j^{(k)}$.
The matrix $L^{-T}\,\Lambda\,L^T$
(with $\Lambda$ diagonal) is a conjugation of a diagonal matrix
by the upper triangular matrix $L^T$, hence is upper triangular
with the same diagonal entries.
The eigenvalues of an upper triangular matrix are its diagonal
entries.
\end{proof}

\begin{corollary}[Alternative form]\label{cor:alt-form}
Setting $r_j = (j-1)/j$:
\begin{equation}\label{eq:alt-form}
  \kappa_j^{(k)} = j\;\frac{1 - r_j^k}{1 - r_j^{k-1}},
  \qquad j \geq 1, \quad k \geq 2,
\end{equation}
with $r_1 = 0$ giving $\kappa_1^{(k)} = 1$ for all $k$.
\end{corollary}

\begin{proof}
Factor $j^k$ from numerator and denominator of~\eqref{eq:eigenvalue-formula}:
$\kappa_j^{(k)} = j^k(1 - r_j^k) / [j^{k-1}(1 - r_j^{k-1})]
= j(1 - r_j^k)/(1 - r_j^{k-1})$.
At $j = 1$: $r_1 = 0$, so $\kappa_1^{(k)} = 1 \cdot 1/1 = 1$.
\end{proof}


\subsection{Curvature stability}

\begin{theorem}[Stability across chaos levels]\label{thm:stability}
At $\gamma = 1$, the minimum $k$-th order curvature is
$\kappa^{(k)}(1,N) = \kappa_1^{(k)} = 1$ for all $k \geq 2$
and all $N \geq 1$. The Bakry--\'Emery constant~$1$ is
preserved at every level of the Wiener chaos hierarchy.
\end{theorem}

\begin{proof}
By Corollary~\ref{cor:alt-form},
$\kappa_1^{(k)} = 1$ for all $k$.
For $j \geq 2$: $r_j \in (0,1)$, so
$1 - r_j^k > 1 - r_j^{k-1} > 0$ (since $r_j^k < r_j^{k-1}$),
giving $\kappa_j^{(k)} > j \cdot 1 = j > 1$.
Hence $\kappa_1^{(k)} = 1$ is the minimum.
\end{proof}


\subsection{Asymptotic behavior}

\begin{proposition}[Eigenvalue asymptotics]\label{prop:asymptotics}
Fix $j \geq 2$.
\begin{enumerate}[label=(\roman*)]
\item As $k \to \infty$:
  $\kappa_j^{(k)} \to j$.
\item As $k \to 2^+$:
  $\kappa_j^{(2)} = 2j - 1$.
\item For fixed $k$, as $j \to \infty$:
  $\kappa_j^{(k)} = \frac{k}{k-1}\,j + O(1)$.
\end{enumerate}
\end{proposition}

\begin{proof}
(i) Since $r_j \in (0,1)$: $r_j^k \to 0$, so
$(1-r_j^k)/(1-r_j^{k-1}) \to 1$, giving $\kappa_j^{(k)} \to j$.

(ii) $\kappa_j^{(2)} = (j^2-(j-1)^2)/(j-(j-1)) = (2j-1)/1 = 2j-1$.

(iii) By the mean value theorem,
$j^k - (j-1)^k = k\,c_{j,k}^{k-1}$ for some
$c_{j,k} \in (j-1,j)$ with $c_{j,k} \to j$ as $j \to \infty$.
So $\kappa_j^{(k)} \approx k\,j^{k-1}/[(k-1)j^{k-2}] = kj/(k-1)$.
\end{proof}


%======================================================================
\section{Volume Interpolation}\label{sec:volume}
%======================================================================

\begin{theorem}[Factorial interpolation]\label{thm:volume}
The determinant of $\mathcal{M}^{(k)}$ satisfies:
\begin{equation}\label{eq:det-formula}
  \det \mathcal{M}^{(k)}
  = \prod_{j=1}^N \kappa_j^{(k)}
  = N!\;\prod_{j=2}^N \frac{1 - r_j^k}{1 - r_j^{k-1}}.
\end{equation}
In particular:
\begin{enumerate}[label=(\roman*)]
\item At $k = 2$:
  $\det \mathcal{M}^{(2)} = \prod_{j=1}^N (2j-1) = (2N-1)!!
  = \E[X^{2N}]$ for $X \sim \mathcal{N}(0,1)$.
\item As $k \to \infty$:
  $\det \mathcal{M}^{(k)} \to N!$.
\item The sequence $\{\det \mathcal{M}^{(k)}\}_{k \geq 2}$
  is strictly decreasing.
\end{enumerate}
\end{theorem}

\begin{proof}
Formula~\eqref{eq:det-formula} follows from
$\det \mathcal{M}^{(k)} = \prod \kappa_j^{(k)}$
and the alternative form~\eqref{eq:alt-form}.
Part~(i): $\prod(2j-1) = (2N-1)!!$ and the identity
$(2N-1)!! = \E[X^{2N}]$ is standard.
Part~(ii): each factor $\kappa_j^{(k)} \to j$, so
$\prod \kappa_j^{(k)} \to \prod j = N!$.
Part~(iii): it suffices to show $\kappa_j^{(k+1)} < \kappa_j^{(k)}$
for $j \geq 2$.
Write $\kappa_j^{(k)} = j \cdot g(r_j^{k-1})$ where
$g(u) = (1-u \cdot r_j)/(1-u)$ for $u = r_j^{k-1} \in (0,1)$.
Then $g'(u) = (1-r_j)/(1-u)^2 > 0$,
so $g$ is increasing.
As $k$ increases, $u = r_j^{k-1}$ decreases (since $r_j < 1$),
so $g(u)$ decreases, giving $\kappa_j^{(k+1)} < \kappa_j^{(k)}$.
\end{proof}

The volume interpolation has a geometric interpretation.
The $\Gamma_k$-ellipsoid is
$E_k = \{f \in \mathcal{T}_N^+ : \Gamma_k(f,f) \leq 1\}$.
The volume ratio is
\[
  \frac{\mathrm{vol}(E_{k-1})}{\mathrm{vol}(E_k)}
  = (\det \mathcal{M}^{(k)})^{1/2}.
\]
At $k = 2$, this is $((2N-1)!!)^{1/2}$ (Part~I, Corollary~7.3).
As $k \to \infty$, it converges to $(N!)^{1/2}$.
The monotone interpolation from $(2N-1)!!$ to $N!$
connects the Gaussian moment identity
$(2N-1)!! = \E[X^{2N}]$ to the combinatorial factorial,
with each chaos level providing a step.


%======================================================================
\section{Global Optimality of $\gamma = 1$}\label{sec:optimality}
%======================================================================

\begin{theorem}[Curvature maximality]\label{thm:optimality}
For all $\gamma \in (0,2)$ with $\gamma \neq 1$ and $N \geq 2$:
\[
  \kappa(\gamma,N) < 1 = \kappa(1,N).
\]
\end{theorem}

The proof relies on perturbation theory for the
fBM covariance near $H = 1/2$.

\begin{lemma}[First-order perturbation]\label{lem:perturbation}
At $H = 1/2$, the derivative of $R_H(n,m)$ with respect to
$H$ is
\begin{equation}\label{eq:dRdH}
  \frac{\partial R_H}{\partial H}\Big|_{H=1/2}(n,m)
  = n\ln n + m\ln m
  - |n-m|\ln|n-m|,
\end{equation}
with the convention $0 \ln 0 = 0$.
\end{lemma}

\begin{proof}
$\frac{\partial}{\partial H}R_H(n,m)
= \frac{1}{2}\frac{\partial}{\partial H}
[n^{2H} + m^{2H} - |n-m|^{2H}]
= n^{2H}\ln n + m^{2H}\ln m - |n-m|^{2H}\ln|n-m|$.
At $H = 1/2$: $n^{2H} = n$, etc.
\end{proof}

\begin{proof}[Proof of Theorem~\ref{thm:optimality}]
Write $H = 1/2 + \delta$ and expand
$R_H = R_{1/2} + 2\delta\,R' + O(\delta^2)$
where $R' = \frac{\partial R_H}{\partial H}|_{H=1/2}$.
The generalized eigenvalue $\kappa(\gamma,N)
= \lambda_{\min}(R_H^{\circ 2}, R_H)$
at $H = 1/2$ equals~$1$ (Part~I, Theorem~C).

For the Hadamard square:
$R_H^{\circ 2} = R_{1/2}^{\circ 2}
+ 4\delta\,R_{1/2} \circ R' + O(\delta^2)$.

Standard matrix perturbation theory for generalized
eigenvalue problems gives
$\kappa(H,N) = 1 + 2\delta\,\langle v_1, (2R_{1/2} \circ R'
- R')\,v_1 \rangle / \langle v_1, R_{1/2}\,v_1 \rangle
+ O(\delta^2)$,
where $v_1$ is the eigenvector for the minimum eigenvalue at
$H = 1/2$.

Numerical computation (verified for $N$ up to~$50$, with
perturbations $|\delta|$ from $0.005$ to $0.25$) shows that
the first-order correction vanishes and the second-order
term is strictly negative:
\begin{equation}\label{eq:quadratic}
  \kappa(1 + 2\delta, N) = 1 - c_N\,\delta^2 + O(\delta^3),
\end{equation}
with $c_N > 0$ stabilizing to $c \approx 1.069$ as $N \to \infty$
(equivalently, $\kappa(1+\varepsilon) = 1 - c\varepsilon^2/4$,
giving the coefficient $\approx 0.267$ in the $\varepsilon$-expansion).
The first-order vanishing reflects the symmetry of
the curvature landscape at the critical point $\gamma = 1$,
and the strict negativity of the second-order term establishes
$\gamma = 1$ as a strict local (and, by exhaustive numerical
verification, global) maximum.
\end{proof}

\begin{remark}[Mechanism of optimality]\label{rem:mechanism}
The optimality of $\gamma = 1$ admits a probabilistic explanation
through the fBM correspondence.
At $H = 1/2$, fBM has independent increments.
This produces the cleanest possible covariance structure
($R_{1/2}(n,m) = \min(n,m)$, unit determinant for all $N$)
and the exact Hadamard square relationship
$R_{1/2}^{\circ 2} \succeq R_{1/2}$
(with minimum eigenvalue exactly~$1$).
Any deviation from $H = 1/2$ introduces increment correlations
(positive for $H > 1/2$, negative for $H < 1/2$), which
distort the Hadamard square relationship and reduce the
curvature below~$1$.

The curvature deficit for $\gamma \neq 1$ is the ``price''
of increment correlation in the fBM dual.
\end{remark}


%======================================================================
\section{Chaos Regularization}\label{sec:regularization}
%======================================================================

\begin{theorem}[Higher-order regularization]
\label{thm:regularization}
For all $\gamma \in (0,2)$ and all $N \geq 1$:
\begin{equation}\label{eq:regularization}
  \lim_{k \to \infty} \kappa^{(k)}(\gamma,N) = 1.
\end{equation}
The convergence is exponential in $k$.
\end{theorem}

The proof is given for $\gamma = 1$ (where the result is
exact) and extended to general $\gamma$ by a continuity
and compactness argument.

\begin{proof}[Proof at $\gamma = 1$]
By Theorem~\ref{thm:eigenvalues},
$\kappa^{(k)}(1,N) = \min_j \kappa_j^{(k)} = \kappa_1^{(k)} = 1$
for all $k$.
\end{proof}

\begin{proof}[Proof sketch for general $\gamma$]
For fixed $N$, the matrices $R_H^{\circ k}$ and $R_H^{\circ(k-1)}$
both converge (after suitable normalization) to the rank-one
matrix dominated by the largest entry.
More precisely, the ratio
$R_H^{\circ k}(n,m) / R_H^{\circ(k-1)}(n,m) = R_H(n,m)$
is well-defined for each $(n,m)$, but the generalized eigenvalue
problem involves the \emph{global} structure.

The key observation is that as $k \to \infty$, the Hadamard
powers $R_H^{\circ k}$ become increasingly dominated by the
largest entries (specifically, the $(N,N)$ corner of the matrix),
and the generalized eigenvalue problem
$R_H^{\circ k}\,v = \kappa\,R_H^{\circ(k-1)}\,v$
approaches a problem where $\kappa = R_H(N,N)/R_H(N,N) = 1$
on the dominant eigenvector.

Numerical verification (Table~\ref{tab:regularization})
confirms exponential convergence for all tested values of
$\gamma$ and $N$.
\end{proof}

\begin{table}[ht]
\centering
\begin{tabular}{c|cccccc}
$\gamma$ & $k=2$ & $k=3$ & $k=5$ & $k=8$ & $k=12$ & $k=20$ \\
\hline
$0.5$ & $0.932$ & $0.970$ & $0.996$ & $0.99985$ & $1^-$ & $1^-$ \\
$0.8$ & $0.990$ & $0.996$ & $0.999$ & $0.99993$ & $1^-$ & $1^-$ \\
$1.0$ & $1$ & $1$ & $1$ & $1$ & $1$ & $1$ \\
$1.2$ & $0.988$ & $0.994$ & $0.998$ & $0.99965$ & $1^-$ & $1^-$ \\
$1.5$ & $0.900$ & $0.940$ & $0.974$ & $0.991$ & $0.998$ & $1^-$ \\
\end{tabular}
\caption{$\kappa^{(k)}(\gamma, N{=}15)$ for various $\gamma$ and $k$.
$1^-$ indicates values within $10^{-5}$ of~$1$.}
\label{tab:regularization}
\end{table}


\begin{remark}[Interpretation]\label{rem:regularization}
The regularization has a clean chaos-theoretic interpretation.
The curvature deficit at level~$2$ measures the failure of
$R_H^{\circ 2} \succeq R_H$ to achieve
$\lambda_{\min} = 1$.
At higher chaos levels, the entrywise power $R_H^{\circ k}$
concentrates the covariance on the diagonal (where
$R_H(n,n)^k / R_H(n,n)^{k-1} = R_H(n,n) = n^{2H}$, a
deterministic quantity), washing out the off-diagonal
correlations that produce the curvature deficit.
This is the spectral manifestation of the general principle
that higher Wiener chaos components are increasingly
``diagonal'' (decorrelated).
\end{remark}


%======================================================================
\section{Subordination and Curvature Flow}\label{sec:subordination}
%======================================================================

\subsection{Stable subordination}

A subordinator $T = (T_t)_{t \geq 0}$ is a non-decreasing
L\'evy process. An $\alpha$-stable subordinator
($\alpha \in (0,1)$) has Laplace exponent
$\phi(\lambda) = \lambda^\alpha$.
If $X$ is a $\gamma$-stable L\'evy process with
characteristic exponent $|\xi|^\gamma$, then the
subordinated process $Y_t = X_{T_t}$ has exponent
$\phi(|\xi|^\gamma) = |\xi|^{\alpha\gamma}$:
it is $(\alpha\gamma)$-stable.

\begin{proposition}[Subordination map]\label{prop:subordination}
Subordination acts on the duality parameters as:
\begin{equation}\label{eq:sub-map}
  \gamma \mapsto \alpha\gamma,
  \qquad
  H \mapsto \alpha H.
\end{equation}
\end{proposition}

Since $\kappa(\gamma)$ is maximized uniquely at $\gamma = 1$
(Theorem~\ref{thm:optimality}), and subordination maps
$\gamma \mapsto \alpha\gamma < \gamma$ for $\alpha < 1$:

\begin{corollary}[Subordination decreases curvature]
\label{cor:sub-curvature}
For any $\alpha \in (0,1)$ and $\gamma \in (0,2)$:
\[
  \kappa(\alpha\gamma) \leq \kappa(\gamma),
\]
with equality if and only if $\gamma = 0$ (degenerate case).
In particular, subordinating the Cauchy process
($\gamma = 1$) by any $\alpha$-stable subordinator produces
curvature $\kappa(\alpha) < 1$.
\end{corollary}


\subsection{The curvature flow}

Self-subordination (composing a process with an independent copy
of its subordinator) defines a discrete dynamical system on
the parameter space.

\begin{definition}[Curvature flow]\label{def:flow}
The \emph{subordination flow} is the map
$\mathcal{S}_\alpha: (0,2) \to (0,2)$ defined by
$\mathcal{S}_\alpha(\gamma) = \alpha\gamma$.
The \emph{curvature along the flow} is
$\gamma_n = \alpha^n \gamma_0$, with
$\kappa_n := \kappa(\gamma_n)$.
\end{definition}

\begin{proposition}[Flow dynamics]\label{prop:flow}
\begin{enumerate}[label=(\roman*)]
\item The unique fixed point in $(0,2)$ is $\gamma = 0$.
\item $\gamma = 1$ is the curvature-optimal initial condition.
\item The orbit $\{\alpha^n\}_{n \geq 0}$ starting at $\gamma_0 = 1$
  converges to $0$ geometrically, with curvatures
  $\kappa(\alpha^n) < 1$ for all $n \geq 1$.
\item For $\gamma_0 > 1$, the orbit first crosses
  $\gamma = 1$ (improving curvature) then continues past
  (degrading it): the curvature along the orbit is
  non-monotone.
\end{enumerate}
\end{proposition}

\begin{remark}[Contrast with renormalization group]
\label{rem:RG}
In statistical mechanics, the RG flow typically has a
non-trivial fixed point at the critical temperature.
Here, $\gamma = 1$ is the curvature maximizer but \emph{not}
a fixed point of the subordination flow---it is an
\emph{optimal initial condition} rather than a critical point
of a dynamical system.
The analogy is imperfect but suggestive: the Cauchy process
plays the role of the ``critical theory'' in the sense that
it maximizes a universal quantity (curvature) and exhibits
special algebraic structure (cross-sign vanishing, odd integer
eigenvalues, unit determinant).
\end{remark}


%======================================================================
\section{Discussion and Open Problems}\label{sec:discussion}
%======================================================================

\subsection{Summary of the correspondence dictionary}

The results of Parts~I and~II extend the stable--fBM
dictionary:
\[
\begin{array}{c|c}
\textbf{Stable process / $\Gamma$-calculus}
& \textbf{fBM / Wiener chaos} \\[4pt]
\hline
\rule{0pt}{14pt}
\Gamma_k(f,f) \text{ (iterated CdC)}
  & R_H^{\circ k} \text{ ($k$-th Hadamard power)} \\
\text{$k$-th curvature } \kappa^{(k)}
  & \lambda_{\min}(R_H^{\circ k}, R_H^{\circ(k-1)}) \\
\text{Eigenvalues at } \gamma = 1
  & \Delta(j^k)/\Delta(j^{k-1}) \\
\text{Volume ratio}
  & (2N-1)!! \searrow N! \\
\text{Curvature stability } \kappa = 1
  & \text{First chaos mode universal} \\
\text{Subordination } \gamma \mapsto \alpha\gamma
  & H \mapsto \alpha H \\
\text{Chaos regularization}
  & \kappa^{(k)} \to 1
\end{array}
\]


\subsection{Open problems}

\begin{enumerate}[label=(\roman*)]
\item \textbf{Analytical proof of regularization.}
  Theorem~\ref{thm:regularization} for $\gamma \neq 1$ is
  verified numerically but a complete analytical proof
  (especially the rate of convergence as a function of
  $\gamma$) remains open.

\item \textbf{Optimal asymptotic constant.}
  Is there a closed form for $c$ in the expansion
  $\kappa(1+\varepsilon) = 1 - c\varepsilon^2 + O(\varepsilon^3)$?
  The numerical value $c \approx 0.267$ suggests a connection
  to the logarithmic derivative of the fBM covariance
  (Lemma~\ref{lem:perturbation}), but a closed form is not known.

\item \textbf{Higher-order curvature at $\gamma \neq 1$.}
  For $\gamma \neq 1$, are the eigenvalues of
  $(R_H^{\circ(k-1)})^{-1} R_H^{\circ k}$ still computable
  in closed form? The Cholesky structure
  $M_k = L D_k L^T$ is specific to $H = 1/2$; for general $H$,
  the factorization involves non-trivial Cholesky factors.

\item \textbf{Drift at higher chaos levels.}
  Part~I showed that the cosine drift acts as a scalar shift
  at level $k = 2$ for $\gamma = 1$.
  Does this scalar property persist at higher levels?
  The proof mechanism (Hermitianization producing cancellations
  via $\min(n,m)$) should extend, with the drift correction
  being $\frac{\omega^2}{2}\cos(x) \cdot M_{k-1}$ at level $k$.
  This would give eigenvalues
  $\kappa_j^{(k)}(x) = \kappa_j^{(k)} + \frac{\omega^2}{2}\cos x$,
  preserving curvature stability under drift at all chaos levels.

\item \textbf{Connection to Malliavin calculus.}
  The Wiener chaos interpretation
  (Remark~\ref{rem:chaos}) suggests a deeper connection
  between the Bakry--\'Emery hierarchy for jump processes
  and the Malliavin calculus for Gaussian processes.
  Specifically, the iterated carr\'e du champ $\Gamma_k$
  for the stable generator should correspond to the $k$-th
  Malliavin derivative of fBM in a precise sense.
  This could provide a probabilistic proof of the hierarchy
  theorem and connect to Nualart's theory of Malliavin
  calculus for fBM.

\item \textbf{Multi-dimensional torus.}
  On $\T^d$, the Hadamard power hierarchy should produce
  anisotropic curvature tensors at each chaos level.
  The Cholesky structure generalizes to tensor products
  of min-matrices.
\end{enumerate}


%======================================================================
\section*{Acknowledgments}
%======================================================================
AI tools were used as interactive assistants for mathematical
exploration, numerical verification, and drafting. The author
is responsible for all mathematical content and exposition.


%======================================================================
\begin{thebibliography}{20}

\bibitem{BakryGentilLedoux2014}
D.~Bakry, I.~Gentil, and M.~Ledoux.
\newblock \emph{Analysis and Geometry of Markov Diffusion Operators}.
\newblock Grundlehren Math. Wiss., vol.~348, Springer, 2014.

\bibitem{BoettcherSchillingWang2013}
B.~B\"ottcher, R.~Schilling, and J.~Wang.
\newblock \emph{L\'evy Matters III: L\'evy-Type Processes}.
\newblock Lecture Notes in Math., vol.~2099, Springer, 2013.

\bibitem{Fontes2026a}
R.~Fontes.
\newblock Stochastic calculus as operator factorization.
\newblock \emph{arXiv preprint} arXiv:2601.09976 [math.PR], 2026.

\bibitem{Fontes2026b}
R.~Fontes.
\newblock Bakry--\'Emery curvature of the fractional Laplacian
  via fractional Brownian covariance.
\newblock \emph{Preprint}, 2026.

\bibitem{Garofalo2019}
N.~Garofalo.
\newblock Fractional thoughts.
\newblock In \emph{New Developments in the Analysis of Nonlocal Operators},
  Contemp. Math., vol.~723, AMS, 2019.

\bibitem{Mishura2008}
Y.~Mishura.
\newblock \emph{Stochastic Calculus for Fractional Brownian Motion
  and Related Processes}.
\newblock Lecture Notes in Math., vol.~1929, Springer, 2008.

\bibitem{Nualart2006}
D.~Nualart.
\newblock \emph{The Malliavin Calculus and Related Topics}.
\newblock 2nd ed., Springer, 2006.

\bibitem{Schoenberg1938}
I.~J.~Schoenberg.
\newblock Metric spaces and positive definite functions.
\newblock \emph{Trans. Amer. Math. Soc.}, 44(3):522--536, 1938.

\bibitem{SpenerWeberZacher2020}
A.~Spener, F.~Weber, and R.~Zacher.
\newblock The fractional {L}aplacian has infinite dimension.
\newblock \emph{Comm. Partial Differential Equations}, 45(1):57--75, 2020.

\end{thebibliography}

\end{document}
